\homework{4}{Wed 21 Feb 2024 18:18}{Exam 1}

\begin{problem}
	Let $A$ be the topological space obtained by deleting $k$ points from the $n$-sphere $S^n$, $n \geq 1, k \geq 1$. Compute $H_*(A)$, i.e, the homology group of $A$ in all non-negative degrees.
\end{problem}
\begin{proof}
	Let $X$ be the space obtained by removing $k$ points from $S^n$. Let $Y$ be the space obtained by removing $k$ points from the interior of $D^n$, so that $X = Y / S^{n - 1}$. Around each removed point in $D^n$ pick an open disk around that point, ensuring the open disks are disjoint, call them $A$. That is, we have a natural embedding $A \hookrightarrow D^n$ such that  $A \cup Y= D^n$.

	Now $A$ and $Y$ form an open cover of $D^n$, so we can apply Mayer-Vietoris sequence of homologies:
% https://quiver.local:8080/#q=WzAsNixbMCwwLCJIX3tufShBIFxcY2FwIFgpIl0sWzIsMCwiSF9uQSBcXG9wbHVzIEhfbiBYIl0sWzQsMCwiSF9uKFNebikiXSxbMCwyLCJIX3tuIC0gMX0oQSBcXGNhcCBYKSJdLFsyLDIsIkhfe24gLSAxfUEgXFxvcGx1cyBIX3tuIC0gMX0gWCJdLFs0LDIsIkhfe24gLSAxfVNebiJdLFswLDEsIlxcYWxwaGFfbiJdLFsxLDIsIlxcYmV0YV9uIl0sWzIsMywiXFxwYXJ0aWFsX24iXSxbMyw0LCJcXGFscGhhX3tuIC0gMX0iXSxbNCw1LCJcXGJldGFfe24gLTF9Il1d
\[\begin{tikzcd}[cramped]
	{H_{n}(A \cap Y)} && {H_nA \oplus H_n Y} && {H_n(D^n)} \\
	\\
	{H_{n - 1}(A \cap Y)} && {H_{n - 1}A \oplus H_{n - 1} Y} && {H_{n - 1}D^n}
	\arrow["{\alpha_n}", from=1-1, to=1-3]
	\arrow["{\beta_n}", from=1-3, to=1-5]
	\arrow["{\partial_n}", from=1-5, to=3-1]
	\arrow["{\alpha_{n - 1}}", from=3-1, to=3-3]
	\arrow["{\beta_{n -1}}", from=3-3, to=3-5]
\end{tikzcd}\]

$A \cap Y$ is disjoint union of $k$ copies of $D_n$ with a point removed, which is homotopic to  $k$ copies of $S^{n - 1}$. Thus the l.e.s. evaluates to
% https://quiver.local:8080/#q=WzAsNixbMCwwLCIwIl0sWzIsMCwiMCBcXG9wbHVzIEhfbiBZIl0sWzQsMCwiMCJdLFswLDIsIlxcWl5rIl0sWzIsMiwiMCBcXG9wbHVzIEhfe24gLSAxfSBZIl0sWzQsMiwiMCJdLFswLDEsIlxcYWxwaGFfbiJdLFsxLDIsIlxcYmV0YV9uIl0sWzIsMywiXFxwYXJ0aWFsX24iXSxbMyw0LCJcXGFscGhhX3tuIC0gMX0iXSxbNCw1LCJcXGJldGFfe24gLTF9Il1d
\[\begin{tikzcd}[cramped]
	0 && {0 \oplus H_n Y} && 0 \\
	\\
	{\Z^k} && {0 \oplus H_{n - 1} Y} && 0
	\arrow["{\alpha_n}", from=1-1, to=1-3]
	\arrow["{\beta_n}", from=1-3, to=1-5]
	\arrow["{\partial_n}", from=1-5, to=3-1]
	\arrow["{\alpha_{n - 1}}", from=3-1, to=3-3]
	\arrow["{\beta_{n -1}}", from=3-3, to=3-5]
\end{tikzcd}\]

This immediately gives $H_{n - 1}Y = \mathbb{Z}^k$. We also have the following diagram for $n = 1$,

% https://quiver.local:8080/#q=WzAsNixbMCwwLCIwIl0sWzIsMCwiMCBcXG9wbHVzIEhfMSBZIl0sWzQsMCwiMCJdLFswLDIsIlxcWiJdLFsyLDIsIjAgXFxvcGx1cyBIX3swfSBZIl0sWzQsMiwiMCJdLFswLDEsIlxcYWxwaGFfbiJdLFsxLDIsIlxcYmV0YV9uIl0sWzIsMywiXFxwYXJ0aWFsX24iXSxbMyw0LCJcXGFscGhhX3tuIC0gMX0iXSxbNCw1LCJcXGJldGFfe24gLTF9Il1d
\[\begin{tikzcd}[cramped]
	0 && {0 \oplus H_1 Y} && 0 \\
	\\
	\Z && {0 \oplus H_{0} Y} && 0
	\arrow["{\alpha_n}", from=1-1, to=1-3]
	\arrow["{\beta_n}", from=1-3, to=1-5]
	\arrow["{\partial_n}", from=1-5, to=3-1]
	\arrow["{\alpha_{n - 1}}", from=3-1, to=3-3]
	\arrow["{\beta_{n -1}}", from=3-3, to=3-5]
\end{tikzcd}\]
So that $H_0 Y = \mathbb{Z}$, and all other homologies of $Y$ are trivial.

Finally we use l.e.s. of a pair to recover $H_*(X)$:
% https://quiver.local:8080/#q=WzAsMTEsWzAsMCwiSF9uKFNee24gLSAxfSkiXSxbMiwwLCJIX24oWSkgPSAwIl0sWzQsMCwiSF9uKFksIFNee24gLSAxfSkiXSxbMCwyLCJIX3tuIC0gMX0oU157biAtIDF9KSA9IFxcWiJdLFsyLDIsIkhfe24gLSAxfShZKSA9IFxcWl5rIl0sWzQsMiwiSF97biAtIDF9KFksIFNee24gLSAxfSkiXSxbMCw0LCJIX3tuIC0gMn0oU157biAtIDF9KSA9IDAiXSxbMiw0LCJIX3tuIC0gMn0oWSkgPSAwIl0sWzQsNCwiSF97biAtMn0oWSwgU157biAtIDF9KSJdLFswLDYsIkhfe24gLSAzfShTXntuIC0gMX0pID0gMCJdLFsyLDYsIlxcY2RvdHMiXSxbMCwxXSxbMSwyXSxbMyw0LCIiLDAseyJjb2xvdXIiOlswLDYwLDYwXSwic3R5bGUiOnsidGFpbCI6eyJuYW1lIjoibW9ubyJ9fX1dLFs0LDVdLFsyLDNdLFs1LDZdLFs2LDddLFs3LDhdLFs4LDldLFs5LDEwXV0=
\[\begin{tikzcd}[cramped]
	{H_n(S^{n - 1})} && {H_n(Y) = 0} && {H_n(Y, S^{n - 1})} \\
	\\
	{H_{n - 1}(S^{n - 1}) = \Z} && {H_{n - 1}(Y) = \Z^k} && {H_{n - 1}(Y, S^{n - 1})} \\
	\\
	{H_{n - 2}(S^{n - 1}) = 0} && {H_{n - 2}(Y) = 0} && {H_{n -2}(Y, S^{n - 1})} \\
	\\
	{H_{n - 3}(S^{n - 1}) = 0} && \cdots
	\arrow[from=1-1, to=1-3]
	\arrow[from=1-3, to=1-5]
	\arrow[color={rgb,255:red,214;green,92;blue,92}, tail, from=3-1, to=3-3]
	\arrow[from=3-3, to=3-5]
	\arrow[from=1-5, to=3-1]
	\arrow[from=3-5, to=5-1]
	\arrow[from=5-1, to=5-3]
	\arrow[from=5-3, to=5-5]
	\arrow[from=5-5, to=7-1]
	\arrow[from=7-1, to=7-3]
\end{tikzcd}\]
The map $H_{n - 1}(S^{n - 1}) \to H_{n - 1}(Y)$ is injective since $Z_{n - 1}(S^{n- 1})$ is freely generated by the $n-1$-simplex $\Delta^{n - 1} \to S^{n - 1}$, and its image in $Y$ is non-contractible. Thus $\partial_n: H_n(Y, S^{n - 1}) \to H_{n - 1}(S^{n - 1})$ is zero. Therefore $H_n(Y, S^{n - 1}) = 0$. Also, the map $H_{n - 1}(Y) \to H_{n - 1}(Y, S^{n - 1})$ is surjective, and we know its kernel is $\mathbb{Z}$, so $H_{n - 1}(Y, S^{n - 1}) \cong \mathbb{Z}^k / \mathbb{Z} = \mathbb{Z}^{k - 1}$.

For the lowest degree, we have the following l.e.s.
% https://quiver.local:8080/#q=WzAsNCxbMiwwLCJIXzAoWSkgPSBcXFoiXSxbNCwwLCJIXzAoWSwgU157biAtIDF9KSJdLFs2LDAsIjAiXSxbMCwwLCJIXzAoU157biAtIDF9KSA9IFxcWiJdLFswLDFdLFsxLDJdLFszLDBdXQ==
\[\begin{tikzcd}[cramped]
	{H_0(S^{n - 1}) = \Z} && {H_0(Y) = \Z} && {H_0(Y, S^{n - 1})} && 0
	\arrow[from=1-3, to=1-5]
	\arrow[from=1-5, to=1-7]
	\arrow[from=1-1, to=1-3]
\end{tikzcd}\]
where, the generator of $H_0(S^{n - 1})$ is indeed a nontrivial $0$-simplex of $H_0(Y)$, but $Y$ is still path-connected, so $H_0(S^{n - 1}) \cong H_0(Y)$. Thus $H_0(Y, S^{n - 1}) = 0$. From those computation we get
\[
	H_m(Y, S^{n - 1}) = \begin{cases}
		\mathbb{Z}^{k - 1} & m = n - 1 \\
		0 & \text{otherwise}
	\end{cases}
.\] 
Finally, consider
\[
	H_*(Y, S^{n - 1}) \cong H_*(X, *) \cong \tilde H_*(X)
.\] 
Hence
\[
	\tilde H_m(X) = \begin{cases}
		\mathbb{Z}^{k - 1} & m = n - 1 \\
		0 & \text{otherwise}
	\end{cases}
.\] 
Thus
\[
		H_m(X) = \begin{cases}
		\mathbb{Z}^{k - 1} & m = n - 1 \\
		\mathbb{Z} & m = 0 \\
		0 & \text{otherwise}
	\end{cases}
.\]
\end{proof}

\begin{problem}
	Let $X$ be any topological space and $I=[0,1]$ be the unit interval. The cone of $X$ is the quotient space $C X:=(X \times I) /(X \times\{0\})$, that is, points at the 'bottom' of $X \times I$ are identified as a single point. $X$ can be embedded to $C X$ by the map $x \mapsto q(x, 1)$ where $q: X \times I \rightarrow C X$ is the quotient map. We identify $X$ with its image in $C X$ and call $X$ a subspace of $C X$. The suspension of $X$ is the quotient space $S X:=C X \cup_X C X$. That is, we take two copies of $C X$, and identify the subspace $X$ in each copy by the identity map $i d: X \rightarrow X$. For instance, the suspension of the $n$-sphere is the $(n+1)$-sphere.\\
a). Show that $C X$ is contractible.\\
b). Let $\tilde{S} X$ be the quotient space of $X \times I$ where we identify $X \times 0$ as a single point and identify $X \times 1$ as a single point. Show $\tilde{S} X$ and $S X$ are homeomorphic.\\
c). Show that for any $n \geq 0, \tilde{H}_{n+1}(S X) \cong \tilde{H}_n(X)$.
\end{problem}
\begin{proof}
	(a) We define a contraction $F$ by the following:

% https://quiver.local:8080/#q=WzAsNCxbMCwwLCJJIFxcdGltZXMgWCBcXHRpbWVzIEkiXSxbMSwwLCJJIFxcdGltZXMgQ1giXSxbMCwxLCJYIFxcdGltZXMgSSJdLFsxLDEsIkNYIl0sWzIsMywicSJdLFswLDEsIlxcdGV4dHtJZH0gXFx0aW1lcyBxIl0sWzAsMiwibSJdLFsxLDMsIkYiXV0=
\[\begin{tikzcd}[cramped]
	{I \times X \times I} & {I \times CX} \\
	{X \times I} & CX
	\arrow["q", from=2-1, to=2-2]
	\arrow["{\text{Id} \times q}", from=1-1, to=1-2]
	\arrow["m", from=1-1, to=2-1]
	\arrow["F", from=1-2, to=2-2]
\end{tikzcd}\]
where $m: (a, x, b) \mapsto (x, ab)$. $m$ sends $I \times X \times \{0\} \to X \times \{0\} $, so $F$ is well-defined and continuous. $F(1, CX) = CX$ and $F(0, CX) = X \times \{0\}  = \{*\} $, so $CX$ is contractible.

(b) Define the homeomorphism $G$ as
% https://quiver.local:8080/#q=WzAsNixbMiwwLCJTWCA9IENYIFxcY3VwX1ggQ1giXSxbNCwxLCJYIFxcdGltZXMgSSJdLFs0LDAsIlxcdGlsZGUgUyBYIl0sWzAsMCwiQ1hfMSBcXHNxY3VwIENYXzIiXSxbMCwxLCIoWCBcXHRpbWVzIEkpXzFcXHNxY3VwKFggXFx0aW1lcyBJKV8yIl0sWzIsMSwiKFggXFx0aW1lcyBJKV8xXFxjdXBfWChYIFxcdGltZXMgSSlfMiJdLFsxLDIsIlxcdGlsZGUgcSJdLFswLDIsIkciXSxbMywwLCJcXHBoaV9YIl0sWzQsMywicSwgcSJdLFs1LDEsImciLDFdLFs0LDUsIlxccHNpX1giLDFdLFs1LDAsInEnLCBxJyJdXQ==
\[\begin{tikzcd}[cramped]
	{CX_1 \sqcup CX_2} && {SX = CX \cup_X CX} && {\tilde S X} \\
	{(X \times I)_1\sqcup(X \times I)_2} && {(X \times I)_1\cup_X(X \times I)_2} && {X \times I}
	\arrow["{\tilde q}", from=2-5, to=1-5]
	\arrow["G", from=1-3, to=1-5]
	\arrow["{\phi_X}", from=1-1, to=1-3]
	\arrow["{q, q}", from=2-1, to=1-1]
	\arrow["g"{description}, from=2-3, to=2-5]
	\arrow["{\psi_X}"{description}, from=2-1, to=2-3]
	\arrow["{q', q'}", from=2-3, to=1-3]
\end{tikzcd}\]
The map $g$ is a homeomorphism: by the following diagram,
% https://quiver.local:8080/#q=WzAsNSxbMCwxLCJYXzIgXFx0aW1lcyBJIl0sWzEsMSwiWCBcXHRpbWVzIEkiXSxbMSwyLCJYXzEgXFx0aW1lcyBJIl0sWzAsMiwiWCJdLFszLDAsIlkiXSxbMCwxLCJcXGlvdGFfMiIsMl0sWzMsMiwiIiwwLHsic3R5bGUiOnsidGFpbCI6eyJuYW1lIjoiaG9vayIsInNpZGUiOiJ0b3AifX19XSxbMywwLCIiLDEseyJzdHlsZSI6eyJ0YWlsIjp7Im5hbWUiOiJob29rIiwic2lkZSI6InRvcCJ9fX1dLFsyLDEsIlxcaW90YV8xIiwyXSxbMiw0XSxbMCw0XSxbMSw0LCIiLDIseyJzdHlsZSI6eyJib2R5Ijp7Im5hbWUiOiJkYXNoZWQifX19XV0=
\[\begin{tikzcd}[cramped]
	&&& Y \\
	{X_2 \times I} & {X \times I} \\
	X & {X_1 \times I}
	\arrow["{\iota_2}"', from=2-1, to=2-2]
	\arrow[hook, from=3-1, to=3-2]
	\arrow[hook, from=3-1, to=2-1]
	\arrow["{\iota_1}"', from=3-2, to=2-2]
	\arrow[from=3-2, to=1-4]
	\arrow[from=2-1, to=1-4]
	\arrow[dashed, from=2-2, to=1-4]
\end{tikzcd}\]
where $\iota_1: (x, a)\mapsto (x, \frac{1}{2} a)$, $\iota_2: (x, a) \mapsto (x, 1 - \frac{1}{2} a)$, we can always construct the dotted arrow given the outer arrows commute. Thus $X \times I$ satisfies the universal property of $(X \times I)_1 \cup _X (X \times I)_2$, so they are the same by definition.
Now we verify that $g$ is compatible with the quotient maps.

% https://quiver.local:8080/#q=WzAsNyxbMCwwLCJDWCJdLFswLDEsIlhfMiBcXHRpbWVzIEkiXSxbMCwyLCJYIl0sWzEsMiwiWF8xIFxcdGltZXMgSSJdLFsxLDEsIlggXFx0aW1lcyBJIl0sWzIsMCwiU1giXSxbMiwyLCJDWCJdLFsxLDAsInFfMiJdLFsyLDEsIiIsMCx7InN0eWxlIjp7InRhaWwiOnsibmFtZSI6Imhvb2siLCJzaWRlIjoidG9wIn19fV0sWzIsMywiIiwyLHsic3R5bGUiOnsidGFpbCI6eyJuYW1lIjoiaG9vayIsInNpZGUiOiJ0b3AifX19XSxbMyw0LCJcXGlvdGFfMSIsMl0sWzEsNCwiXFxpb3RhXzIiLDJdLFswLDVdLFszLDYsInFfMSJdLFs2LDVdLFs0LDUsIihxJywgcScpIiwxXSxbMiwwLCIiLDEseyJjdXJ2ZSI6LTIsInN0eWxlIjp7InRhaWwiOnsibmFtZSI6Imhvb2siLCJzaWRlIjoidG9wIn19fV0sWzIsNiwiIiwxLHsiY3VydmUiOjIsInN0eWxlIjp7InRhaWwiOnsibmFtZSI6Imhvb2siLCJzaWRlIjoidG9wIn19fV1d
\[\begin{tikzcd}[cramped]
	CX && SX \\
	{X_2 \times I} & {X \times I} \\
	X & {X_1 \times I} & CX
	\arrow["{q_2}", from=2-1, to=1-1]
	\arrow[hook, from=3-1, to=2-1]
	\arrow[hook, from=3-1, to=3-2]
	\arrow["{\iota_1}"', from=3-2, to=2-2]
	\arrow["{\iota_2}"', from=2-1, to=2-2]
	\arrow[from=1-1, to=1-3]
	\arrow["{q_1}", from=3-2, to=3-3]
	\arrow[from=3-3, to=1-3]
	\arrow["{(q', q')}"{description}, from=2-2, to=1-3]
	\arrow[curve={height=-12pt}, hook, from=3-1, to=1-1]
	\arrow[curve={height=12pt}, hook, from=3-1, to=3-3]
\end{tikzcd}\]
In the above diagram the arrow $(q', q')$ uniquely exists. $(q', q')$ can be defined by going around the commutative diagram for $X \times [0,\frac{1}{2}] $ and $X \times [\frac{1}{2}, 1]$ respectively, and they agree at $X \times \{\frac{1}{2}\} $.

The pair of quotients $q$ gives rise to a quotient map $(q', q')$ that makes the diagram commute; we also have another quotient map $\tilde q$. To show $SX \cong \tilde S X$, we need to show $(q', q') X \times I \cong \tilde q X \times I$. What is $(q', q')$ doing to $X \times I$? Well, since $q_2$ sends $X_2 \times \{0\}$ to the point $[X_2 \times \{0\} ]$ in $CX$,  $(q', q')$ has to send $\iota_2(X_2 \times \{0\} ) = X \times \{1\} $ to a point in $SX$. Similarly, $(q', q')$ collapses $\iota_1(X_1 \times \{0\} ) = X \times \{0\} $ to a point in $SX$. From this we know that $(q', q')$ is exactly same as the map $\tilde q$ which identifies $X \times \{0\} $ and $X \times \{1\} $. Hence $SX \cong \tilde SX$.
%To see $SX \cong \tilde SX$, we only need to see $\operatorname{ker} \tilde q = \operatorname{ker} (q', q')$. We know  $\operatorname{ker} (q', q') = (\iota_1, \iota_2)(\operatorname{ker} q_1, \operatorname{ker} q_2) = X \times \{0\} \sqcup X \times \{1\}  $. But $\operatorname{ker} \tilde q = X \times \{0\} \sqcup X \times \{1\}  $. So those two constructions are the same.
\end{proof}



\begin{problem}
	3. CW complexes are constructed by gluing cells. A question is how the CW complex depends on the attaching maps. Here we address this question.\\
In this problem, a deformation retraction always means a strong deformation retraction. Let $(X, A)$ be a pair of topological spaces such that $A \subset X$ and $X \times I$ deformation retracts to $X \times\{0\} \cup A \times I$. Let $Y$ be a topological space, and $f: A \times I \rightarrow Y$ be a continuous map. We use the notation $f_t(a):=f(a, t), a \in A, t \in I$. So $f$ is a homotopy between $f_0$ and $f_1$.\\
a). Prove $X \times I$ also deformation retracts to $X \times\{1\} \cup A \times I$.\\
b). For $i=0,1$, denote by $Y \cup_{f_i} X$ be the space obtained by gluing $X$ to $Y$ along the map $f_i: A \rightarrow Y$, that is, $Y \cup_{f_i} X$ is the quotient space $Y \sqcup X / \sim$, where the equivalence relation is $a \sim f_i(a), a \in A$. Similarly, define the space $Y \cup_f(X \times I)$. Prove that $Y \cup_f(X \times I)$ deformation retracts both to $Y \cup_{f_0} X$ and to $Y \cup_{f_1} X$. Therefore, $Y \cup_{f_0} X$ and $Y \cup_{f_1} X$ are homotopy equivalent.\\
c). You should convince yourself or take it as a fact that $\left(D^n, S^{n-1}\right)$ satisfies the condition of the pair $(X, A)$ described in the body of the problem. Hence, the homotopy class of a CW complex depends only on the homotopy class of its attaching maps. Consider the 'dunce hat', as shown Figure 1, which is obtained by identifying the three edges of a triangle (with interior) according to the arrows shown in the figure.\\
It is a CW-complex with one 0-cell, one 1-cell, and one 2-cell. The 1-skeleton is hence a circle and is homeomorphic to the edge labelled by $a$. Consider $a$ as an oriented loop. The attaching map for the 2-cell, as a loop in the 1-skeleton, is given by $a * a * \bar{a}$, where $\bar{a}$ is the loop $a$ with reversed orientation, and ' $*$ ' means concatenation of loops. Prove that the dunce hat is contractible.
\end{problem}
\begin{proof}
	(a) Let $F$ be the deformation retract $I \times (X \times I) \to (X \times I)$ such that $F(1, X \times I) = X \times \{0\}  \cup A \times I$.  Define the ``flip map'' $\phi: X \times I \to X \times I$ by $(x, a) \mapsto (x, 1 - a)$. Now $\phi$ is a homeomorphism. The map $(1 \times \phi) F (1 \times \phi)^{-1}$ is the required deformation retract.

	(b) By symmetry given in (a), we only need to show deformation retract to $Y \cup _{f_0} X$. Let $F: I \times (X \times I) \to X \times I$ be the deformation retract $X \times I \curvearrowright X \times \{0\}  \cup A \times I$. We can trivially extend it to $F: I \times \left( Y \sqcup X \times I \right) \to Y \sqcup (X \times I)$ by being identity on $Y$ ; this turns it into a deformation retract $Y \sqcup (X \times I) \curvearrowright Y \sqcup (X \times \{0\}  \cup A \times I)$.

	Consider the following diagram:

% https://quiver.local:8080/#q=WzAsNCxbMCwwLCJJIFxcdGltZXMgWSBcXGN1cF9mIChYIFxcdGltZXMgSSkiXSxbMCwyLCJJIFxcdGltZXMgWSBcXHNxY3VwIChYIFxcdGltZXMgSSkiXSxbMiwyLCJZIFxcc3FjdXAgKFggXFx0aW1lcyBJKSJdLFsyLDAsIlkgXFxjdXBfZiAoWCBcXHRpbWVzIEkpIl0sWzEsMCwiKDEscSkiXSxbMSwyLCJGIiwyXSxbMiwzLCJxIiwyXSxbMCwzLCJcXHRpbGRlIEYiXV0=
\[\begin{tikzcd}[cramped]
	{I \times Y \cup_f (X \times I)} && {Y \cup_f (X \times I)} \\
	\\
	{I \times Y \sqcup (X \times I)} && {Y \sqcup (X \times I)}
	\arrow["{(1,q)}", from=3-1, to=1-1]
	\arrow["F"', from=3-1, to=3-3]
	\arrow["q"', from=3-3, to=1-3]
	\arrow["{\tilde F}", from=1-1, to=1-3]
\end{tikzcd}\]

	We verify that $F$ maps through the quotients. The map $(1,  q)$ identifies all pairs  $\{(t,(a, s)), (t,f(a, s))\} $ for $(a, s) \in A \times I$. Now $F(t, (a,s)) = (a, s)$ since $F$ is a strong deformation retract. $f(t, f(a, s)) = f(a, s)$ since $f(a, s) \in Y$. Moreover, $\{(a, s), f(a, s)\} $ is identified by the quotient map $q$.

	Thus, $F$ factors through the quotient to give a continuous map
	\[
		\tilde F: I \times (Y \cup _f (X \times I)) \to Y \cup _f (X \times I)
	.\] 
	We verify that this is a deformation retract to $Y \cup_{f_0} X$.
	\begin{itemize}
		\item  $\tilde F\left( 0, Y \cup _f (X \times I) \right) = q F\left( 0, Y \sqcup \left( X \times I \right)  \right)  = q \left( Y \sqcup (X \times I) \right) = Y \cup _f (X \times I)$.
		\item Similarly, $\tilde F(1, Y \cup _f (X \times I)) = q \left( Y \sqcup X \times \{0\}  \right)$. Since $f_0(a) = f(a, 0)$, the middle square of the following diagram commutes (dotted line represent natural isomorphism):
% https://quiver.local:8080/#q=WzAsOCxbMiwwLCJZIl0sWzIsMSwiQSBcXHRpbWVzIFxcezBcXH0iXSxbMywxLCJZIFxcc3FjdXAgWCBcXHRpbWVzIFxcezBcXH0iXSxbMywwLCJZIFxcY3VwX2YgWCBcXHRpbWVzIFxcezBcXH0iXSxbMCwwLCJZIFxcY3VwX3tmXzB9IFgiXSxbMSwwLCJZIl0sWzEsMSwiQSJdLFswLDEsIlkgXFxzcWN1cCBYIl0sWzEsMCwiZiJdLFsxLDIsIiIsMix7InN0eWxlIjp7InRhaWwiOnsibmFtZSI6Imhvb2siLCJzaWRlIjoidG9wIn19fV0sWzAsM10sWzIsMywicSIsMl0sWzUsNF0sWzYsNSwiZl8wIl0sWzcsNCwicSIsMl0sWzYsNywiIiwyLHsic3R5bGUiOnsidGFpbCI6eyJuYW1lIjoiaG9vayIsInNpZGUiOiJ0b3AifX19XSxbNSwwLCIiLDEseyJzdHlsZSI6eyJib2R5Ijp7Im5hbWUiOiJkYXNoZWQifSwiaGVhZCI6eyJuYW1lIjoibm9uZSJ9fX1dLFs2LDEsIiIsMSx7InN0eWxlIjp7ImJvZHkiOnsibmFtZSI6ImRhc2hlZCJ9LCJoZWFkIjp7Im5hbWUiOiJub25lIn19fV0sWzcsMiwiIiwxLHsib2Zmc2V0IjozLCJjdXJ2ZSI6Miwic3R5bGUiOnsiYm9keSI6eyJuYW1lIjoiZGFzaGVkIn0sImhlYWQiOnsibmFtZSI6Im5vbmUifX19XV0=
\[\begin{tikzcd}[cramped]
	{Y \cup_{f_0} X} & Y & Y & {Y \cup_f X \times \{0\}} \\
	{Y \sqcup X} & A & {A \times \{0\}} & {Y \sqcup X \times \{0\}}
	\arrow["f", from=2-3, to=1-3]
	\arrow[hook, from=2-3, to=2-4]
	\arrow[from=1-3, to=1-4]
	\arrow["q"', from=2-4, to=1-4]
	\arrow[from=1-2, to=1-1]
	\arrow["{f_0}", from=2-2, to=1-2]
	\arrow["q"', from=2-1, to=1-1]
	\arrow[hook, from=2-2, to=2-1]
	\arrow[dashed, no head, from=1-2, to=1-3]
	\arrow[dashed, no head, from=2-2, to=2-3]
	\arrow[shift right=3, curve={height=12pt}, dashed, no head, from=2-1, to=2-4]
\end{tikzcd}\]
, thus by universal property, there is a unique isomorphism between $Y \cup _f X \times \{0\} $ and $Y \cup_{f_0} X$.

\item Finally we verify that $\tilde F$ fixes $Y \cup_{f_0} X$ at all times, i.e. the deformation retraction is strong. Take $y \in Y$. Then $\tilde F(t, q(y)) = \tilde F((1, q)(t, y)) = q(F(t, y)) = q(y)$, so $\tilde F$ fixes $q(Y)$. Next consider $(a, s) \in A \times I$. Then $\tilde F(t, q(a, s)) = q\left( F(t, (a, s)) \right)  = q(a, s)$, so $\tilde F$ fixes $q(A \times I )$. Finally, consider $(x, 0) \in X \times \{0\} $. Then  $\tilde F(t, q(x, 0)) = q\left( F(t, (x, 0)) \right)  = q(x, 0)$. Thus, $F$ restricted to $I \times \left( Y \sqcup \left( X \times \{0\}  \cup A \times I \right)  \right) $ is a projection onto $Y \sqcup \left( X \times \{0\}  \cup A \times I \right)$, and this factors through the quotient to give $\tilde F$ fixing  $q\left(Y \sqcup \left( X \times \{0\}  \cup A \times I \right)    \right)  = Y \cup _f (X \times \{0\}) \cong Y \cup _{f_0} X$.
\end{itemize}

(c) Let $\gamma$ denote the generator of $1$st homology of $S^{1}$, and let $f_0$ be the map $\gamma \mapsto a$, $f_1$ be the map $\gamma \mapsto a a \overline{a} $. Here $a$ can be identified viewed as a path $a: [0,1] \to Y$. Define $a_t: [0,1] \to Y$ to be the path defined by $a_t(s) = a(ts)$. Then $a_1 = a$ and $a_0 = \text{Id}$. Now define the homotopy by
\[
	f(t, s) = a * a_t * \overline{a_t} (s)
.\] 
By construction this is a continuous map. Also $f(0, s) = a* a_0 * \overline{a_0}  = a$ and  $f(1, s) = a * a* \overline{a} $.

Now that we have the homotopy, part (2) tells us that
\[
	D^2 \cup _{f_1} S^1 \simeq  D^2 \cup _{f_0} S^1 \cong D^2 \simeq  *
.\] 
where $\simeq $ means homotopic and $\cong$ means homeomorphic.
\end{proof}

\begin{problem}
	4. Let $X$ be a topological space, and $p \in S^1$ be a point in $S^1$. Consider the space $X \times S^1$.\\
a). Denote by $i: X \times\{p\} \rightarrow X \times S^1$ the embedding. Prove that $i$ induces injective maps on homology, that is, for any $n \geq 0, i_*: H_n(X \times\{p\}) \rightarrow H_n\left(X \times S^1\right)$ is injective.\\
b). For any $n \geq 0$, prove that $H_n\left(X \times S^1\right) \cong H_n(X) \oplus H_n\left(X \times S^1, X \times\{p\}\right)$.\\
c). For any $n \geq 0$, prove that $H_n\left(X \times S^1, X \times\{p\}\right) \cong H_{n-1}(X)$. Therefore, we have $H_n\left(X \times S^1\right) \cong H_n(X) \oplus H_{n-1}(X)$.
\end{problem}
\begin{proof}
	(a) From the long exact sequence
% https://quiver.local:8080/#q=WzAsNSxbMiwwLCJIX3tuICsgMX0oWCBcXHRpbWVzIFNeMSkiXSxbNCwwLCJIX3tuICsgMX0oWCBcXHRpbWVzIFNeMSwgWCBcXHRpbWVzIFxce3BcXH0pIl0sWzAsMSwiSF8gbihYKSJdLFsyLDEsIkhfe259KFggXFx0aW1lcyBTXjEpIl0sWzQsMSwiSF97bn0oWCBcXHRpbWVzIFNeMSwgWCBcXHRpbWVzIFxce3BcXH0pIl0sWzAsMV0sWzMsNCwial9uIl0sWzIsMywiaV9uIl0sWzEsMiwiXFxwYXJ0aWFsIiwyXV0=
\[\begin{tikzcd}[cramped]
	&& {H_{n + 1}(X \times S^1)} && {H_{n + 1}(X \times S^1, X \times \{p\})} \\
	{H_ n(X)} && {H_{n}(X \times S^1)} && {H_{n}(X \times S^1, X \times \{p\})}
	\arrow[from=1-3, to=1-5]
	\arrow["{j_n}", from=2-3, to=2-5]
	\arrow["{i_n}", from=2-1, to=2-3]
	\arrow["\partial"', from=1-5, to=2-1]
\end{tikzcd}\]
, and the fact that the map $H_{n + 1}(X \times S^1) \to H_{n + 1}(X \times S^1, X \times \{p\} )$ is a surjection, we know that the map $\partial$ is zero, hence $i_n$ is injective.

(b) The same argument justifies that all boundary maps are zero. Hence we can use the split lemma to the short exact sequence $0 \to H_n(X) \to H_n(X \times S^1) \to H_n(X \times S^1, X \times \{p\} ) \to 0$. We only need to verify that the s.e.s. splits. To do so, we explicitly construct a map $r: H_n(X \times S^1) \to H_n(X \times \{p\} )$. We do so by direct projection. The projection $S_n(X \times S^1) \to S_n(X \times \{p\} )$ is a surjective chain map. Also, $r i_n = \text{Id}$. This verifies condition for split lemma.

(c) Use Relative Mayler-Vietoris Sequence (Hatcher, p.152)

\begin{lemma}
	(Hatcher p.152)
	A relative form of the Mayer-Vietoris sequence is sometimes useful. If one has a pair of spaces $(X, Y)=(A \cup B, C \cup D)$ with $C \subset A$ and $D \subset B$, such that $X$ is the union of the interiors of $A$ and $B$, and $Y$ is the union of the interiors of $C$ and $D$, then there is a relative Mayer-Vietoris sequence
\[
\cdots \rightarrow H_n(A \cap B, C \cap D) \xrightarrow{\Phi} H_n(A, C) \oplus H_n(B, D) \xrightarrow{\Psi} H_n(X, Y) \xrightarrow{\partial} \cdots
\]
\end{lemma}

Now, write  
\begin{itemize}
	\item $A = X \times (S^1 \setminus \{q\}) $, $B = X \times (S^1 \setminus \{r\} )$, $C = D = X \times  \{p\} $, where $p, q, r$ all distinct. Then all assumptions of the lemma are satisfied. 

		To write down the sequence we observe that $H_n(A, C) = H_n(B, D) = 0$ because $X \times S^1 \setminus \{q\} $ is homotopic to $X \times \{p\} $. This directly gives for all $n$,
		\[
			H_{n}(X \times S^1, X \times \{p\} ) \cong H_{n - 1} (A \cap B, C \cap D)
		.\] 
		But $X \times \{S^1 \setminus \{q, r\} \} \simeq X \times \{p\}  \sqcup X \times \{p\} $, so
		\[
			H_{n - 1}(A \cap B, C \cap D) = H_{n - 1}(X \times \{p\} ) = H_{n - 1}(X)
		.\] 
\end{itemize}

\end{proof}



\begin{problem}
	5. Let $\Sigma_g$ be the closed oriented surface of genus $g$. Embed it in $\mathbb{R}^3$ in a 'standard' way so that it is the boundary of a handlebody $U_g$. (So $U_1$ is just a solid torus.) On $\Sigma_g$, choose a set of closed curves $m_1, \cdots, m_g$ and and $l_1, \cdots, l_g$. These $2 g$ curves are mutually disjoint except that each $m_i$ intersects $l_i$ exactly once. Each $m_i$ bounds a disk in $U_g$. Choose an orientation on each curve such that $m_i$ and $l_i$ intersect positively, that is, at the intersection, when you rotate the arrow on $m_i$ counterclockwise by 90 degree, it will match the arrow on $l_i$. See Figure 2. Treat each closed curve as a 1-cycle. Then $H_1\left(\Sigma_g\right)$ is freely generated by $\left\{\left[m_1\right], \cdots,\left[m_g\right],\left[l_1\right], \cdots,\left[l_g\right]\right\}$.\\
a). Show that $H_1\left(U_g\right)$ is freely generated by $\left\{\left[l_1\right], \cdots,\left[l_g\right]\right\}$.\\
b). If $\gamma$ is an oriented simple closed curve on $\Sigma_g$, give a combinatorial way (e.g., in terms of counting intersections) to express $[\gamma]$ as a linear combination of the $\left[m_i\right]^{\prime} \mathrm{s}$ and $\left[l_i\right]^{\prime} s$ in $H_1\left(\Sigma_g\right)$. You don't have to explain why if you are confident of your answer.\\
c). Let $\phi: \Sigma_g \rightarrow \Sigma_g$ be a homeomorphism, then $\phi_*: H_1\left(\Sigma_g\right) \rightarrow H_1\left(\Sigma_g\right)$ is an isomorphism. Let\\
\[
\phi_*\left(\left[m_j\right]\right)=\sum_{i=1}^g A_{i j}\left[m_i\right]+\sum_{j=1}^g B_{i j}\left[l_i\right], \quad A_{i j}, B_{i j} \in \mathbb{Z} .
\]

In practice, given $\phi$, we can use the method in part $a$ ) to find the coefficients $A_{i j}, B_{i j}$. Let $B=\left(B_{i j}\right)$ be the $g$ by $g$ matrix and think of it as $\mathbb{Z}$-linear map on $\mathbb{Z}^g$ sending a column vector $v \in \mathbb{Z}^n$ to $B v$. Let $M_\phi$ be the space obtained by gluing two copies of $U_g$ along $\phi$. That is, let $U_g^1:=U_g$ and $U_g^2:=U_g$, and treat $\phi$ as a $\operatorname{map} \phi: \partial U_g^1 \rightarrow \partial U_g^2$, then
\[
M_\phi:=\left(U_g^2 \cup_\phi U_g^1\right) .
\]

Prove that $H_2\left(M_\phi\right) \cong \operatorname{ker}(B)$. In particular, it means $H_2\left(M_\phi\right)$ is free Abelian.\\
The space $M_\phi$ is in fact a closed oriented 3-manifold. The decomposition $M_\phi=U_g \cup_\phi U_g$ is called a Heegaard decomposition of $M_\phi$.
\end{problem}
\begin{proof}
	(a) Use induction on $g$. $U_g$ is obtained by attaching $U_{1}$ to $U_{g - 1}$. Apply reduced Mayer-Vietoris, we have
% https://quiver.local:8080/#q=WzAsOCxbMCwwLCJcXHRpbGRlIEhfMShBIFxcY2FwIEIpIl0sWzIsMCwiXFx0aWxkZSBIXzEoQSkgXFxvcGx1cyBcXHRpbGRlIEhfMShCKSJdLFs0LDAsIlxcdGlsZGUgSF8xKFVfZykiXSxbNiwwLCJcXHRpbGRlIEhfMChBIFxcY2FwIEIpIl0sWzAsMSwiMCJdLFsyLDEsIlxcWlxcbGFuZ2xlW2xfMV0sIFxcbGRvdHMsW2xfe2cgLSAxfV1cXHJhbmdsZSBcXG9wbHVzIFxcWlxcbGFuZ2xlW2xfZ11cXHJhbmdsZSJdLFs0LDEsIlxcdGlsZGUgSF8xKFVfZykiXSxbNiwxLCIwIl0sWzAsMV0sWzEsMl0sWzIsM10sWzQsNV0sWzUsNl0sWzYsN11d
\[\begin{tikzcd}[cramped]
	{\tilde H_1(A \cap B)} && {\tilde H_1(A) \oplus \tilde H_1(B)} && {\tilde H_1(U_g)} && {\tilde H_0(A \cap B)} \\
	0 && {\Z\langle[l_1], \ldots,[l_{g - 1}]\rangle \oplus \Z\langle[l_g]\rangle} && {\tilde H_1(U_g)} && 0
	\arrow[from=1-1, to=1-3]
	\arrow[from=1-3, to=1-5]
	\arrow[from=1-5, to=1-7]
	\arrow[from=2-1, to=2-3]
	\arrow[from=2-3, to=2-5]
	\arrow[from=2-5, to=2-7]
\end{tikzcd}\]
Combine with the fact that reduced homology coincides with homology for $n > 0$, we obtain the claim.

(b) The coefficient for $[l_k]$ is the number of times $\gamma$ intersects $m_k$ positively (in the sense that their cross product points outward the surface), minus the number of times $\gamma$ intersects $m_k$ negatively.

(c) 
\end{proof}

